\section{Smooth lattice polytopes and projective charts}
\label{sec:delzant-projective-charts}
The slopes determine a compactification as well as a growth condition. At a smooth vertex the primitive edge directions are a lattice basis, so the characters in those directions become ordinary holomorphic coordinates across the boundary.

\Needspace{8\baselineskip}
\begin{definition}\label{def:delzant-lattice-polytope}
A lattice polytope in $M_{\mathbb R}$ is the convex hull of finitely many lattice points. A face is the nonempty set on which some linear functional has its minimum; the polytope itself is allowed. Vertices are zero-dimensional faces, edges are one-dimensional faces, and facets have codimension one. A full-dimensional lattice polytope is smooth, or Delzant, if at each vertex its primitive edge directions form a basis of $M$. In dimension zero the single point is smooth, with the empty basis.
\end{definition}
\index{polytope!lattice}\index{polytope!Delzant}

The cone generated by a set consists of its finite nonnegative linear combinations. A ray of a cone is extreme if a sum of two cone vectors belonging to that ray has both summands on the ray.

\begin{lemma}\label{lem:polytope-facets-and-vertex-cones}
A full-dimensional lattice polytope has an irredundant facet presentation
\[
 P=\{m:\ell_\rho(m)=\langle m,u_\rho\rangle+a_\rho\geq0
                                      \text{ for every facet }\rho\},
\]
where $u_\rho\in N$ is primitive and inward pointing and $a_\rho\in\mathbb Z$. At a vertex $v$ its tangent cone is
\[
 C_v=\{y:\langle y,u_\rho\rangle\geq0
                                      \text{ for facets containing }v\}
     =\operatorname{cone}(P-v).
\]
Its extreme rays are the edge directions at $v$. If $P$ is smooth, exactly $n$ facets contain $v$; their normals $u_1,\ldots,u_n$ are dual to the primitive edge basis $e_1,\ldots,e_n$. Thus
\[
 m-v=\sum_i t_i e_i,\qquad t_i=\ell_i(m)\geq0\quad(m\in P),
\]
and $v+e_i\in P\cap M$.
\end{lemma}
\begin{proof}
Discard a generating point whenever it belongs to the convex hull of the others. The remaining points are exactly the vertices: a nontrivial segment expression would express that point by the other generators after moving its own coefficient to the left, while strict separation from their convex hull exposes the point as a zero-dimensional face. In particular the vertices are integral. Each exposed face is the convex hull of the vertices minimizing its functional, so there are finitely many faces.

Every boundary point lies in a facet. Indeed a supporting hyperplane gives a proper exposed face containing it. If that face has dimension below $n-1$, vary its supporting normal within the annihilator of the face's directions, in a direction independent of the original normal. Choose the sign so that the value at some outside vertex decreases relative to the face. Increase the variation until the first additional vertex reaches the supporting value. There are finitely many outside vertices, all with positive initial gaps, so this first parameter exists and is positive. The normal remains nonzero by independence. The enlarged face has greater affine dimension. Repeating reaches a facet. A facet contains $n$ affinely independent integral vertices; Gaussian elimination gives a rational normal, hence its unique primitive integral inward normal, and its constant is integral. For a point outside $P$, travel to it from an interior point. The first boundary point lies in a facet; continuing the same line violates that facet inequality. Thus the facet half-spaces intersect exactly in $P$. At a relative interior point of a facet all other inequalities are strict; a small outward move violates only this one, proving irredundancy.

At $v$ the active inequalities give the displayed cone. Every vector in that cone gives a point $v+ty\in P$ for all sufficiently small $t>0$, since the finitely many inactive inequalities are strictly positive at $v$. This proves the tangent-cone identity. The cone is pointed: a line in it would give a segment of $P$ with $v$ in its interior, contradicting extremality. The sum of its active normals is strictly positive on every nonzero cone vector; otherwise all active pairings vanish and the opposite vector also lies in the cone. Slicing by value one of that sum therefore gives a bounded polytope.

A bounded intersection of finitely many half-spaces is the convex hull of its vertices. In its affine span a point is a vertex exactly when its active equalities have full rank. Otherwise choose a nonzero direction annihilating them. Boundedness permits extension in both signs to boundary points with an additional independent active equality. Express the original point as their convex combination and repeat. The rank strictly increases until vertices are reached. Applied to the cone slice, this proves generation by extreme rays. At a slice vertex, the active cone equalities cut out a one-dimensional ray; the same equalities in $P$ cut out an edge beginning at $v$. Conversely the direction of an exposed edge must be extreme, since a decomposition into two cone directions would, for small $t$, express a point of that edge by two points of $P$ outside its exposed face.

For a smooth polytope these rays are the basis vectors $e_i$, so the cone is the corresponding orthant. Its facets have the primitive dual normals. Every facet of $P$ through $v$ gives a cone facet, and every cone facet gives the corresponding local facet of $P$; hence there are exactly $n$. The nonnegative coordinate formula follows. The other endpoint of the edge in direction $e_i$ is integral, so its length in that primitive direction is a positive integer. It therefore contains $v+e_i$. The dimension-zero assertions have no facets or edges.
\end{proof}

Fix a smooth $P$ for the rest of the chapter. Put $A=P\cap M$ and $T=\operatorname{Hom}(M,\mathbb C^*)$, with compact subgroup $T_c=\operatorname{Hom}(M,S^1)$. A lattice point $m$ defines the character $\chi^m(t)=t(m)$; choosing a lattice basis identifies $T$ with $(\mathbb C^*)^n$. Define
\[
 \iota:T\longrightarrow\mathbb P^{|A|-1},\qquad
                       t\longmapsto[\chi^m(t)]_{m\in A},
 \qquad X_P=\overline{\iota(T)}.
\]
The closure here is in ordinary complex projective space: its points are nonzero vectors modulo scalar multiplication, and $Z_v\ne0$ has coordinates $Z_m/Z_v$. This space is compact as the quotient of the unit sphere by scalar phases. The map $[Z]\mapsto ZZ^*/|Z|^2$ identifies it with the compact set of rank-one Hermitian orthogonal projectors, so it is Hausdorff. The action $\tau\cdot[Z_m]=[\chi^m(\tau)Z_m]$ preserves $X_P$.

\begin{theorem}\label{thm:delzant-projective-monomial-charts}
The space $X_P$ is a smooth projective variety of dimension $n$, and $\iota$ identifies $T$ with a dense open subset. For each vertex $v$, the set $U_v=\{Z_v\ne0\}\cap X_P$ is an affine chart isomorphic to $\mathbb C^n$, with coordinates $w_i=Z_{v+e_i}/Z_v=\chi^{e_i}$ on $T$.
These vertex charts cover $X_P$. Writing $m-v=\sum_i b_i e_i$ gives $Z_m/Z_v=\prod_i w_i^{b_i}$ with nonnegative integers $b_i$.
The restricted hyperplane bundle $L=\mathcal O_{\mathbb P^{|A|-1}}(1)|_{X_P}$ is very ample and has a smooth positive quadratic metric.
\end{theorem}
\begin{proof}
For each $m\in A$, choose a convex combination of affinely independent vertices equal to $m$, using \cref{lem:finite-dimensional-convex-separation}. Its coefficients are rational by Gaussian elimination on the integral coordinates. Clearing denominators gives
\[
 d m=\sum_v c_v v,\qquad c_v\in\mathbb Z_{\geq0},\quad\sum_v c_v=d,
 \qquad Z_m^d=\prod_v Z_v^{c_v}
\]
on the torus and hence on its closure. Thus at least one vertex coordinate is nonzero at every point of the closure.

The preceding lemma gives nonnegative integral coordinates $b_i$ for $m-v$. On $Z_v\ne0$, the torus therefore lies in the graph of the displayed monomials in $w_i$. Conversely that entire graph belongs to $X_P$, since $(\mathbb C^*)^n$ is dense in $\mathbb C^n$. Its closure in this affine chart is the graph. This proves the chart assertion, smoothness and dimension; the coordinates $w_i$ recover every torus point, proving injectivity and the open embedding.

It also proves algebraicity globally. In addition to the finitely many vertex-combination equations above, impose, for each vertex and each $m\ne v$, the homogeneous equations
\[
 Z_m Z_v^{\sum_i b_i-1}=\prod_i Z_{v+e_i}^{b_i}.
\]
Their common projective zero set has a nonzero vertex coordinate by the first equations, and on that chart is exactly the monomial graph by the second. It is therefore $X_P$. Two proper algebraic closed subsets cannot cover it: a defining polynomial for each can be chosen with a nonzero Laurent restriction to $T$, whereas their product would vanish there. Laurent polynomials form a domain, giving the contradiction. Thus $X_P$ is irreducible.

The coordinate sections of $L$ have no common zero and give this projective embedding, which is the definition of very ampleness. On $U_v$ give its nonvanishing coordinate section $s_v$ squared norm $\exp(-H_v)$, where
\[
 H_v=\log\sum_{m\in A}|w^{b(m)}|^2.
\]
These are the restricted Fubini--Study weights and have the required metric transition law. Their Hessians are positive definite: for a tangent vector, Cauchy--Schwarz makes the numerator of the Hessian of the logarithmic sum nonnegative. Equality requires the derivatives of all monomials to be proportional to their values. The constant monomial forces the proportionality factor to vanish, and the monomials $w_i$ then force the tangent vector to vanish. Curvature is $\ddc H_v$ with the book's normalization, so this is a smooth positive metric. In dimension zero $X_P$ is a point and positivity has its usual vacuous meaning.
\end{proof}
\index{toric variety!projective monomial charts}

\section{Fans, divisors and character sections}
\label{sec:toric-fans-and-character-sections}
For a face $Q$ let $\sigma_Q$ be the cone generated by the inward normals of facets containing $Q$; put $\sigma_P=\{0\}$. The following identifies the boundary independently of a general toric-variety construction.

\begin{proposition}\label{prop:normal-fan-and-uniform-ray-coefficients}
The cones $\sigma_Q$ form a complete smooth fan: they cover $N_{\mathbb R}$, each is generated by part of a lattice basis, all its faces occur in the family, and any two cones intersect in a face of each. Equivalently $x\in\sigma_Q$ if the minimum of $\langle m,x\rangle$ on $P$ is attained on all of $Q$.
For $n>0$ there is $\Lambda$ depending only on $P$ and the fixed Euclidean norm such that every unit $x$ has an expression
\[
 x=\lambda\sum_i c_i u_{\rho_i},\qquad
 0<\lambda\leq\Lambda,\quad c_i>0,\quad\sum_i c_i=1,
\]
with all rays in a single cone of this fan.
\end{proposition}
\begin{proof}
Choose a vertex $v\in Q$. If $v$ minimizes the pairing, evaluating along each initial edge gives $x=\sum_i\alpha_i u_i$ with $\alpha_i=\langle e_i,x\rangle\geq0$. If all of $Q$ minimizes it, the formula $m-v=\sum_i t_i e_i$ shows that $\alpha_i=0$ whenever the corresponding facet does not contain $Q$. Conversely a nonnegative combination of normals of facets containing $Q$ plainly has its minimum there. This proves the characterization and the lattice-basis assertion. Selecting a subset of these generators corresponds to intersecting their facets at $v$, giving the normal cone of that larger face. Thus every cone face occurs. Every functional has a minimum face, giving completeness. The least face containing two given faces is the intersection of all facets containing both. The characterization therefore says that the two normal cones intersect in its normal cone; its generating rays are a subset of those of each original cone, so this is a face of each.

For a unit vector choose a minimizing vertex. Its nonnegative coefficients satisfy $\alpha_i\leq|e_i|$. Their sum is positive and bounded by $\sum_i|e_i|$. Take $\Lambda$ to be the maximum of these finite sums over all vertices, divide by the sum, and omit zero coefficients. In dimension zero the fan consists of its zero cone and there is no unit vector.
\end{proof}
\index{fan!inner normal fan}

\begin{proposition}\label{prop:toric-orbits-and-prime-divisors}
The torus orbits in $X_P$ are indexed by the nonempty faces $Q$ of $P$ and have dimension $\dim Q$. The orbit closure for $Q$ contains exactly the orbits indexed by its faces. The complement of $T$ is the union of the invariant prime divisors $D_\rho$ indexed by facets. For every character,
\[
 \operatorname{ord}_{D_\rho}\chi^m=\langle m,u_\rho\rangle.
\]
The monomial charts construct the toric variety of the inner normal fan just described.
\end{proposition}
\begin{proof}
In a vertex chart, the torus scales the $w_i$ independently. An orbit is therefore specified by the set $I$ of vanishing coordinates, with all others nonzero. Associate to it the face $Q$ cut out by $\ell_i=0$ for $i\in I$. It has dimension $n-|I|$: near $v$, the remaining independent edge directions occur with small positive coefficients. A projective coordinate $Z_m$ is nonzero on this orbit exactly when $m\in Q\cap A$. This criterion is independent of the chart; the convex hull of those lattice points is $Q$, so different faces give different orbits. The coordinates outside $Q$ stay zero on its orbit closure, excluding every face not contained in $Q$. For a face $Q'\subset Q$, choose a vertex of $Q'$ and use its chart. Let the additional coordinates defining $Q'$ tend to zero while keeping the others nonzero. This approximates its orbit by points of the orbit of $Q$, proving the converse inclusion.

A facet orbit has codimension one, and its closure is an irreducible divisor: its dense orbit is a complex torus, whose Laurent coordinate ring has no zero divisors. In a vertex chart meeting that orbit it is $w_i=0$, with primitive normal $u_i$. Since $\chi^m=\prod_j w_j^{\langle m,u_j\rangle}$ on the torus, its order is the claimed integer. Finally the semigroup of characters nonnegative on a vertex cone is generated freely by its dual edge basis. Thus its affine character chart is exactly the $\mathbb C^n$ chart already constructed, and the common character identifications give their gluing. This is the fan construction in the present smooth projective case.
\end{proof}
\index{toric variety!orbit--face correspondence}\index{divisor!toric prime}

\begin{theorem}\label{thm:toric-character-section-basis}
There is a meromorphic section $s_0$ of $L$ with divisor
\[
 D_P=\sum_\rho a_\rho D_\rho,
\]
so that $L=\mathcal O_{X_P}(D_P)$. For every integer $k\geq0$,
\[
 H^0(X_P,L^k)=\bigoplus_{m\in kP\cap M}\mathbb C\,\chi^m s_0^k.
\]
Here $0P=\{0\}$. Each summand is a character line for the natural torus action, and
\[
 \lim_{k\to\infty}\frac{n!}{k^n}h^0(X_P,L^k)=n!\operatorname{vol}(P).
\]
\end{theorem}
\begin{proof}
On the open torus the coordinate sections satisfy $s_v/s_w=\chi^{v-w}$. Hence $s_0=s_v/\chi^v$ is independent of the vertex. In a vertex chart its coefficient in the frame $s_v$ is a Laurent monomial, so it defines a global meromorphic section. The order formula gives $\operatorname{ord}_{D_\rho}s_0=-\langle v,u_\rho\rangle=a_\rho$ when $v$ belongs to that facet. There are no other zeros or poles because $s_0$ is nonvanishing on $T$. The local meromorphic equations of this divisor are $\chi^{-v}$; their inverses are the local frames of $\mathcal O(D_P)$, with transition $\chi^{v-w}$, exactly the transition of $L$.

Write a global section on $T$ as $f s_0^k$. Iterating the one-variable Laurent theorem gives a unique expansion $f=\sum_{m\in M}c_m\chi^m$, normally convergent on compact polyannuli. In $U_v$ its coefficient in the frame $s_v^k$ is $f\chi^{-kv}$. Holomorphy on $\mathbb C^n$ forces all Laurent terms with a negative exponent in an edge coordinate to vanish, by uniqueness of its Taylor and Laurent coefficients. Thus $c_m\ne0$ implies $\langle m-kv,u_i\rangle\geq0$ for every active facet at every vertex. Every facet contains a vertex, so these are exactly the inequalities $m\in kP$ for $k>0$. For $k=0$ a nonzero allowed vector would give $p+tm\in P$ for all $t\geq0$, contradicting boundedness. There are finitely many allowed points. Conversely for such an $m$ all its local exponents are nonnegative, so $\chi^m s_0^k$ extends in every chart. Uniqueness of Laurent coefficients proves the direct sum. The torus action preserves each monomial line, giving the character lines. The last limit is \cref{thm:convex-body-lattice-riemann-sums}, with lattice volume one and the displayed factor $n!$. Dimension zero gives one section and volume one.
\end{proof}
\index{line bundle!toric character basis}

\section{The Guillemin reference metric}
\label{sec:guillemin-reference-metric}
The projective construction supplies a positive metric, but the logarithmic behavior of a specific reference metric can matter in an integrability calculation. The following entropy potential has exactly the prescribed facet logarithms.

\begin{theorem}\label{thm:guillemin-reference-metric}
Set $0\log0=0$ and define the closed convex function
\[
 G_0(m)=\sum_\rho\ell_\rho(m)\log\ell_\rho(m)\quad(m\in P),
 \qquad G_0=+\infty\quad\text{outside }P.
\]
Put $F_0=G_0^*$. Then $F_0$ is smooth on $N_{\mathbb R}$, has positive definite Hessian, and $F_0\sim h_P$. The gradient $\nabla G_0$ is a diffeomorphism from $\operatorname{Int}P$ onto $N_{\mathbb R}$. With $\operatorname{Trop}(t)(m)=\log|\chi^m(t)|^2$, the local weights
\[
 H_v=F_0\circ\operatorname{Trop}-\log|\chi^v|^2
\]
extend smoothly from $T$ to $U_v$ and define a smooth positive quadratic metric on $L$, invariant under $T_c$. Its curvature is the unique smooth extension of $\ddc(F_0\circ\operatorname{Trop})$ and represents $c_1(L)$.
\end{theorem}
\begin{proof}
The scalar function $t\log t$ is convex and continuous on $[0,+\infty)$. Thus $G_0$ is closed and convex, continuous and bounded on $P$. Its Hessian on the interior is
\[
 D^2G_0=\sum_\rho\frac{u_\rho\otimes u_\rho}{\ell_\rho},
 \qquad
 \nabla G_0=\sum_\rho(\log\ell_\rho+1)u_\rho.
\]
The normals span, so the Hessian is positive definite. For each $x$, minimize $G_0(m)-\langle m,x\rangle$ on $P$. A boundary minimizer is impossible: moving toward a fixed interior point gives terms $t b\log(t b)$ for its vanishing facets, whose right derivative is $-\infty$; all other terms have finite derivatives. The minimum is therefore interior and unique. The inverse function theorem and the positive Hessian show that the gradient is a global diffeomorphism. Legendre differentiation gives $\nabla F_0=(\nabla G_0)^{-1}$ and $D^2F_0=(D^2G_0)^{-1}$. Boundedness of $G_0$ on $P$ gives $F_0\sim h_P$.

Fix a vertex chart and a point with $w_i=0$ for $i\in I$, and $w_j\ne0$ for $j\in J$, where $I,J$ partition the edge indices. Write $t_i=\ell_i(m)$ for its vertex-facet coordinates, so $m=v+\sum_i t_i e_i$. The corresponding face $Q$ has $t_i=0$ for $i\in I$. For the prescribed real numbers $x_j=\log|w_j|^2$, $j\in J$, minimizing $G_0|_Q-\sum_{j\in J}t_jx_j$ gives a unique point $q$ in the relative interior of $Q$. The preceding boundary-derivative argument works inside that face as well. Strict positivity there follows already from the terms $t_j\log t_j$ for $j\in J$.

Near $q$ all other facet functions are positive, and
\[
 G_0(t)=\sum_{i\in I}t_i\log t_i+g(t)
\]
with $g$ smooth. The interior gradient equations are
\[
 x_i=\log t_i+1+g_i(t)\quad(i\in I),\qquad
 x_j=g_j(t)\quad(j\in J).
\]
Replace the first equations by $y_i=t_i\exp(1+g_i(t))$, with $y_i=|w_i|^2$. At $q$ the derivative of the map $t\mapsto(y_I,x_J)$ has a positive diagonal $I$ block, a zero upper-right block and the positive definite block $(g_{jk})_{j,k\in J}$ in the lower right. It is invertible. The real inverse function theorem therefore solves $t$ smoothly as a function of $(y_I,x_J)$ near this boundary point. Moreover
\[
 t_i=y_i\exp(-1-g_i(t))\quad(i\in I),
\]
so the solution lies in $P$ for nonnegative sufficiently small $y_I$, with its other facets still positive.

On the open torus the local weight, after subtracting $\langle v,x\rangle$, is
\[
 H_v=\sum_{i\in I}t_i(1+g_i(t))
                  +\sum_{j\in J}t_jx_j-g(t).
\]
The singular logarithms have canceled. This expression is smooth in $(|w_i|^2)_{i\in I}$ and $(\log|w_j|^2)_{j\in J}$, proving smooth extension near the chosen point. To check strict positivity there, Legendre differentiation gives $\partial H_v/\partial x_i=t_i$. Hence at $y_i=0$,
\[
 \frac{\partial H_v}{\partial y_i}
                       =\exp(-1-g_i(q))>0.
\]
The complex Hessian has these positive numbers as its diagonal $I$ block; all mixed entries involving a vanishing coordinate and a different coordinate vanish. In the nonvanishing $J$ block the Hessian in logarithmic coordinates is $(g_{JJ})^{-1}$, hence positive definite, and passing to $w_J$ is an invertible diagonal congruence. This proves positivity at the boundary, as well as on the open torus.

On chart overlaps $H_v-H_w=\log|\chi^{w-v}|^2$, which is the logarithm of the squared absolute value of a holomorphic unit. Thus $|s_v|^2=\exp(-H_v)$ satisfies the quadratic metric transition law. The weights depend only on character moduli, and $T_c$ acts on the frames by unit factors, proving metric invariance. The curvatures glue, have class $c_1(L)$, and agree with the stated form on $T$. Density gives uniqueness of the smooth extension. In dimension zero the facet sum and both potentials are zero, $L$ is trivial on a point and its curvature is zero.
\end{proof}
\index{Guillemin metric}\index{Legendre transform!Guillemin potential}
